% MLP+CP quantiles (surrogate/raw, n_cal=400):
% alpha=0.10 -> {'log10_P': 1.852863464608622, 'log10_K': 0.7210642050153828, 'e': 0.6437431975955525, 'omega': 3.043812346536477}
% alpha=0.40 -> {'log10_P': 1.2216077948944926, 'log10_K': 0.3886337770549111, 'e': 0.506152803441515, 'omega': 2.718299843987451}
% coverage block: {'per_coord_coverage': {'log10_P': 1.0, 'log10_K': 0.9824561403508771, 'e': 1.0, 'omega': 0.9649122807017544}, 'per_coord_median_width': {'log10_P': 2.602156315870957, 'log10_K': 1.2040900820736484, 'e': 0.8854764774220711, 'omega': 6.111998329451626}, 'joint_coverage': 0.9473684210526315, 'frac_infinite': 0.0}
\subsection{Exoplanet Radial-Velocity Experiment}
\label{sec:exoplanet}

We apply the same simulation-based conformal pipeline to single-planet Radial Velocity (RV) time series from the NASA Exoplanet Archive.
The latent parameters are
\[
\theta = (\log_{10} P,\; \log_{10} K,\; e,\; \cos\omega,\; \sin\omega),
\]
where $P$ is the orbital period, $K$ the RV semi-amplitude, $e$ the eccentricity, and $\omega$ the argument of periastron.
The noiseless forward map $h(\theta)$ is the Keplerian RV model evaluated on the observation times; synthetic curves are generated by sampling $\theta$ from empirical priors fitted on tabulated real systems, resampling observation cadences, and adding residual noise from a trained sparse variational GP.

\paragraph{Point predictor.}
The regressor $\psi$ is a multilayer perceptron trained on synthetic curves with a 74-dimensional summary+LSP feature vector and a dual eccentricity head (separate circular / eccentric pathways gated by an $e>0$ classifier).
On a held-out synthetic validation split it recovers $(\log_{10} P, \log_{10} K, e)$ as shown in Figure~\ref{fig:rv-pred}.
The 74-D feature set does not encode periapsis epoch, so $\omega$ is weakly identifiable from these features alone; we therefore omit $\omega$ from the published pred-vs-true panel and treat $\omega$ conformal intervals as largely uninformative (near full-circle half-widths).
A stable alternative targetisation uses Laplace coordinates $(k,h)=(e\cos\omega, e\sin\omega)$ (already available as \texttt{--targets hk} in our codebase); recovering $e=z_2/\cos\omega$ from $(\omega, e\cos\omega)$ is singular at $\omega=\pm\pi/2$ and is not recommended.

\paragraph{Uncertainty.}
Prediction regions come from the simulation-based conformal procedure of this paper (split conformal calibrated on synthetic curves, evaluated under synthetic$\to$real covariate shift), with conformity scores on the physical coordinates $(\log_{10} P, \log_{10} K, e, \omega)$ and Bonferroni aggregation.
We use the MLP as $\psi$ with the surrogate-label score.
On the real test split (likelihood-ratio reweighted), joint coverage at $1-\alpha=0.90$ is approximately 0.95, with $\log_{10}P$ median width $2.60$, $\log_{10}K$ $1.20$, $e$ $0.89$, $\omega$ $6.11$~rad (median support-capped widths).
These regions are intentionally conservative under covariate shift with a moderately accurate $\psi$; $\alpha{=}0.1$ half-widths are reported in the Earth-like summary table and $\alpha{=}0.4$ regions are used for the light phase-folded overlays.

\paragraph{Figures.}
Figure~\ref{fig:rv-heldout} shows a held-out real RV series phase-folded at the tabulated period, with (i) observations, (ii) $h(\theta_{\mathrm{tab}})$, (iii) $h(\psi(y))$, (iv) light traces for $\theta\sim\mathrm{Unif}(\Gamma_{0.4})$, and (v) noisy simulator draws at $\psi(y)$.
Figure~\ref{fig:rv-pred} shows synthetic validation pred-vs-true scatters for $\log_{10} P$, $\log_{10} K$, and $e$.

\begin{figure}[t]
\centering
\includegraphics[width=\linewidth]{rv_heldout_phasefold.png}
\caption{Held-out real RV (HD~139357) phase-folded at the tabulated period. Overlays: observations; noiseless Keplerian from tabulated and predicted parameters; $\sim$20 light traces from the $\alpha{=}0.4$ conformal region; noisy simulator draws at $\psi(y)$.}
\label{fig:rv-heldout}
\end{figure}

\begin{figure}[t]
\centering
\includegraphics[width=\linewidth]{rv_pred_vs_true.png}
\caption{Synthetic validation pred-vs-true for the 74-D MLP ($P$, $K$, $e$). Argument of periastron is omitted because the 74-D features lack periapsis-epoch information.}
\label{fig:rv-pred}
\end{figure}
